\documentclass[10pt,a4paper]{amsart}
\usepackage[margin=19mm]{geometry}
\usepackage{amsmath,amssymb,mathtools}
\usepackage[scr=rsfs,cal=euler]{mathalfa}
\usepackage{xfrac}
\usepackage[unicode,hidelinks,bookmarks=false]{hyperref}
\newcommand{\ie}{i.e.\ }
\newcommand{\cf}{cf.\ }
\newcommand{\resp}{respectively\ }
\newcommand{\Subsection}{\subsection}
\providecommand{\nobreakdash}{\nobreak\hbox{-}}
\setlength{\emergencystretch}{2em}
\multlinegap0pt
\usepackage{bm}
\usepackage{bbm}
\usepackage{dsfont}
\usepackage{xspace}
\usepackage{cancel}
\usepackage{mathtools}
\usepackage{amsthm}
\makeatletter
\@ifundefined{thm}{\newtheorem{thm}{Thm}}{}
\@ifundefined{lem}{\newtheorem{lem}[thm]{Lemma}}{}
\@ifundefined{prop}{\newtheorem{prop}[thm]{Proposition}}{}
\makeatother
\newcommand{\mc}[1]{\mathcal{#1}}
\newcommand{\ol}[1]{\overline{#1}}
\newcommand{\op}{\textup{op}}
\title[On the convergence of explicit formulas for $L^2$ solutions ]{On the convergence of explicit formulas for $L^2$ solutions to the Benjamin-Ono and continuum Calogero-Moser equations}
\author{Yvonne Alama Bronsard, Thierry Laurens}
\date{Source: arXiv:2602.19046v1 (CC BY 4.0)}
\begin{document}
\maketitle

\noindent\emph{Source and licence.} On the convergence of explicit formulas for $L^2$ solutions to the Benjamin-Ono and continuum Calogero-Moser equations. Version arXiv:2602.19046v1 (CC BY 4.0), distributed under the Creative Commons Attribution 4.0 International licence, as recorded in the arXiv licence metadata for this exact version and shown on the abstract page https://arxiv.org/abs/2602.19046v1. Full rights notice: licenses/arxiv-2602.19046-CC-BY-4.0.txt.\par\smallskip

\noindent\textit{Editorial scope.} Counted unit: the Prop whose source label is \texttt{lem:resolv} in the article (source0.tex lines 910--916, source label \texttt{lem:resolv}), delivered together with the article's own complete original proof of it (source0.tex lines 918--982, one \texttt{proof} environment of 65 lines). No mathematics is written, repaired or re-derived: statement and proof are the article's own text line for line. Required context retained uncounted: eq:BO (lines 119--121); eq:CS (lines 135--137); prop:lax (lines 715--733); R (lines 726--731); prop:lax' (lines 765--779); R' (lines 772--777); lem:resolv 2 (lines 801--804); lem:resolv 2' (lines 860--863). Every definition, display and label of those lines is retained unchanged. Omitted from this package and not redistributed: 1--118, 122--134, 138--714, 732--764, 778--800, 805--859, 864--909, 917--917, 983--1154. Nothing inside a retained excerpt is omitted or altered: every retained body is delivered exactly as the article's lines are written, so no omission receipt of any kind accompanies this package. Citations: the retained excerpts carry no citation command at all, so no bibliography entry of the article is transcribed and no reference is redistributed. Reference adaptation: none was needed and none was made; every reference inside the delivered unit resolves inside it, so no unresolved reference or citation is produced. Printed slips kept literal: no printed word, symbol, hypothesis or number of the counted statement or of its proof was altered. Layout and class: the article's own class is \texttt{article} with packages including authblk, amsrefs, subcaption, graphicx, amssymb, amsmath, placeins, mathrsfs, multirow, amsthm, float, bbm, colortbl, xcolor; the wrapper is the lane's pinned \texttt{amsart} editorial wrapper at 10pt on A4, which restates the theorem-like environments the retained text needs and adds the article's own macro definitions that the retained text needs (\texttt{\textbackslash mc}, \texttt{\textbackslash ol}, \texttt{\textbackslash op}), with extra wrapper packages (bm, bbm, dsfont, xspace, cancel, mathtools). The delivered numbering is the unit's own. Assets: no retained span contains a figure environment, a graphics-inclusion command, a PDF-page inclusion, a table or a TikZ picture; no image file is copied into this package, the package declares no asset dependency and no image file is redistributed with it. Licence: version arXiv:2602.19046v1 of this article is distributed under the Creative Commons Attribution 4.0 International licence, as recorded in the arXiv licence metadata for this exact version and shown on the abstract page https://arxiv.org/abs/2602.19046v1. A full rights notice is delivered at licenses/arxiv-2602.19046-CC-BY-4.0.txt. Characters: every retained excerpt is pure ASCII, so the delivered engine, XeLaTeX with its default Latin Modern text font, reports no missing character.
\begin{equation}\label{eq:BO}\tag{BO}
\partial_t u(t,x) = \partial_x \left( |\partial_x| u - u^2 \right)(t,x), \quad  u_{|t=0}(x) = u_0(x), \quad (t,x) \in \mathbb{R} \times \mathbb{T}, 
\end{equation}

\begin{equation}\label{eq:CS}\tag{CCM}
i\partial_t u + \partial_x^2u \pm \frac{2}{i} u\partial_x \Pi(|u|^2) = 0,\quad  u_{|t=0}(x)= u_{0}(x), \quad (t,x) \in \mathbb{R} \times \mathbb{T}. 
\end{equation}%

\begin{prop}[The truncated Lax operators for \eqref{eq:BO}]
\label{prop:lax}
Given $u\in L^2$ real-valued and $1 \le n \le \infty$, the operator
\begin{equation*}
L_n  = -i\partial_x - \Pi_n  u \Pi_n 
\end{equation*}
with domain $D(L_n) = H^1_+$ is self-adjoint and semi-bounded.  Moreover, there is a constant $\kappa_0 = \kappa_0(\|u\|_{L^2}) \geq 1$ such that for all $\kappa\geq\kappa_0$, the resolvent of $L_n$,
\[
R_n(\kappa;u) = (L_n + \kappa)^{-1}
\] 
exists, and maps $L^2_+$ into $H^{1}_+$. Moreover, the following estimates hold 
\begin{align}
\tfrac12\|f\|_{H^1_\kappa} &\le \| (L_n+\kappa)f \|_{L^2} \le \tfrac32 \|f\|_{H^1_\kappa} ,
\label{L} \\
\tfrac23\| f\|_{H^{-1}_\kappa} &\le\;\; \| R_n(\kappa)f \|_{L^2}\;\; \le2 \| f\|_{H^{-1}_\kappa}
\label{R}
\end{align}
uniformly for $1\leq n\leq \infty$ and $\kappa\geq\kappa_0$.
\end{prop}

\begin{prop}[The truncated Lax operators for \eqref{eq:CS}]
\label{prop:lax'}
Given $u\in L^2_+$ and $1\leq n\leq \infty$, the operator
\begin{equation*}
L_n f = -i\partial f \mp \Pi_n [ u \Pi_n (\ol{u} \Pi_nf) ]
\end{equation*}
with domain $D(L_n) = H^1_+$ is self-adjoint and semi-bounded.  Moreover, there is a constant $\kappa_0 = \kappa_0(u) \geq 1$ so that for all $\kappa\geq\kappa_0$, the resolvent $R_n(\kappa;u)$ of $L_n$ exists, maps $L^2_+$ into $H^{1}_+$, and 
\begin{align}
\tfrac12\|f\|_{H^1_\kappa} &\le \| (L_n+\kappa)f \|_{L^2} \le \tfrac32 \|f\|_{H^1_\kappa} ,
\label{L'} \\
\tfrac23\| f\|_{H^{-1}_\kappa} &\le\;\; \| R_n(\kappa)f \|_{L^2}\;\; \le2 \| f\|_{H^{-1}_\kappa}
\label{R'}
\end{align}
uniformly for $1\leq n\leq \infty$.
\end{prop}

\begin{lem}[Norm resolvent convergence for \eqref{eq:BO}]
\label{lem:resolv 2}
Given a real-valued function $u\in L^2$, the operators $L_n$ converge to $L_\infty$ in norm resolvent sense as $n\to\infty$.
\end{lem}

\begin{lem}[Norm resolvent convergence for \eqref{eq:CS}]
\label{lem:resolv 2'}
Given $u\in L^2_+$, the operators $L_n$ converge to $L_\infty$ in norm resolvent sense as $n\to\infty$.
\end{lem}

\begin{prop}\label{lem:resolv}
For $1\leq n\leq \infty$, let $L_n$ denote the truncated Lax operators for \eqref{eq:BO} or \eqref{eq:CS} constructed in Proposition~\ref{prop:lax} or~\ref{prop:lax'} for any fixed initial data $u$.  Then, for any $T>0$ and any compact set $\mc{F}\subset L^2_{+}$, we have
\begin{equation}
\sup_{f\in \mathcal{F}}\ \sup_{t\in[-T,T]}\, \| (e^{itL_n} - e^{itL_\infty})f \|_{L^2} \xrightarrow{n \to \infty} 0.
\label{resolv}
\end{equation}
\end{prop}

\begin{proof}
We will present the argument in the settings of~\eqref{eq:BO} and~\eqref{eq:CS} in parallel, indicating the source of each ingredient for both cases when necessary.

Set $\kappa = \kappa_0$, where $\kappa_0 = \kappa_0(u)$ is defined in Proposition~\ref{prop:lax} and~\ref{prop:lax'} for~\eqref{eq:BO} and~\eqref{eq:CS} respectively.
For $M \ge 1$, let $\chi_{\leq M}: \mathbb{R} \to  [0,1]$ be a smooth function that is equal to one on $[-\kappa, M]$, and goes to zero as $x\to +\infty$. We also let $X = [-T,T] \times [ -\kappa, +\infty)$ and define
%
\[ g \colon
\begin{array}{rccl}
 &X & \longrightarrow & \mathbb{R} \\
&(t, x) & \longmapsto & e^{itx} \chi_{\leq M}(x).
\end{array}\]
%

By the Stone--Weierstrass theorem, polynomials in $t$ and $(x+\kappa)^{-1}$ are dense in $C_{0}(X)$.  Therefore, for any $\epsilon>0$ there exists a polynomial $P(t,y)$ 
such that
\[
\sup_{(t,x)\in X} \left|P\big(t,(x+\kappa)^{-1}\big) - g(t,x)\right| \le \epsilon.
\] 
Hence, by the functional calculus we have
\begin{equation}\label{eq:secondT}
\sup_{\substack{t \in [-T,T] \\[.2em] 1\le n \le \infty }} \left\| P\big(t, (L_n + \kappa)^{-1}\big) - g(t, L_n)\right\|_{\op} \le \epsilon.
\end{equation}
%

We can then decompose the error into three terms,
\begin{align}\label{eq:3terms}
\| (e^{itL_n} - e^{itL_\infty})f \|_{L^2}\nonumber
{}\leq{}& 2\sup_{\substack{t \in [-T,T] \\[.2em] 1\le n \le \infty }}\| e^{itL_{n}}f - g(t, L_n)f \|_{L^2} \\
&+ 2C\sup_{\substack{t \in [-T,T] \\[.2em] 1\le n \le \infty }} \big\|  P\big(t, (L_n+\kappa)^{-1}\big) - g(t, L_n) \big\|_{\op} \\\nonumber
&+ C\sup_{t\in[-T,T]}\big\|  P\big(t, (L_n+\kappa)^{-1}\big) - P\big(t,(L+\kappa)^{-1}\big) \big\|_{\op},
\end{align}
with $C = \sup_{f\in \mathcal{F}} \| f\|_{L^2}$. The second term in \eqref{eq:3terms} can be taken arbitrarily small using \eqref{eq:secondT}. We now prove that the third term in \eqref{eq:3terms} goes to zero as $n\to \infty$. Indeed, for $m\in \mathbb{N}$ we have
\begin{align*}
R_n^{m} - R_\infty^{m} = R_n^{m-1} &\left( R_n - R_\infty \right) 
+\left(R_n^{m-1} - R_\infty^{m-1} \right)R_\infty,
\end{align*}
and hence, using \eqref{R} and \eqref{R'},
\[
\|R_n^{m} - R^{m}_\infty \|_\op \le 2^{m-1}\|R_n - R_\infty \|_\op + 2
\|R_n^{m-1} - R_\infty^{m-1}\|_\op.
\]
By induction, $\|R_n^{m} - R^{m}_\infty \|_\op \le 3^m \|R_n - R_\infty \|_\op$, which goes to zero as $n\to \infty$ by Lemmas~\ref{lem:resolv 2} and \ref{lem:resolv 2'}.
Next, by writing the polynomial as $P(t,y) = \sum_{\ell,m}\alpha_{\ell,m} t^{\ell} y^{m}$, it follows that
\[
\sup_{t\in[-T,T]}\|  P\big(t, (L_n+\kappa)^{-1}\big) - P\big(t,(L+\kappa)^{-1}\big)\|_{\op} \le \sum_{\ell, m} \alpha_{\ell,m} T^{\ell} \|R_n^{m} - R_\infty^{m} \|_\op \xrightarrow{n\to\infty} 0 .
\]
%
Finally, we show that the first term in \eqref{eq:3terms} goes to zero as $M \to \infty$, finishing the proof. As $\mathcal{F} \subset L^2_+$ is compact and $H^1_{+}$ is dense in $L^2_{+}$, we have that given any $\epsilon>0$ there exists $m\in\mathbb{N}$ and $\phi_1, \dots, \phi_m \in H^1_{+}$ such that 
\[
\forall f\in\mathcal{F}, \ \exists j \in \{1,\dots, m\},\ \|\phi_j - f\|_{L^2} \le \epsilon.
\]
Thus, by letting $\chi_{>M} = 1-\chi_{\le M}$ we have that for $t\in[-T,T]$ and $n\in\mathbb{N}$
\[
\| e^{itL_{n}}f - g(t, L_n)f \|_{L^2} \le \| \chi_{> M}(L_n)e^{itL_n}f\|_{L^2} \le \| \chi_{>M}(L_n) (f-\phi_i) \|_{L^2} + \| \chi_{>M}(L_{n})\phi_i\|_{L^2}.
\]
The first term in the right-hand side above is bounded by $\epsilon$ while for the second term, using the functional calculus and inequalities \eqref{R} and \eqref{R'}, we have
\[
\| \chi_{>M}(L_n) \phi_i\|_{L^2} \le \left\| \chi_{>M}(L_n)\frac{L_n+\kappa}{M+\kappa}\phi_i\right\|_{L^2}
\le \frac{1}{M} \| (L_n+ \kappa)\phi_i\|_{L^2}
\le \frac{3}{2M} \max_{1\le i\le m}\| \phi_i\|_{H^1_{\kappa}},
\]
which is independent of $n$ and goes to zero as $M\to \infty$.
To conclude the proof we take first $M$, then $\deg(P)$, and finally $n$ large enough, in order for the three terms in \eqref{eq:3terms} to be arbitrarily small. \qedhere

\end{proof}

\end{document}
